\documentclass[11pt]{article}

\usepackage{amsmath,amssymb}
\usepackage{xurl}
\usepackage[hidelinks]{hyperref}
\setlength{\emergencystretch}{2em}

% Shared commands for notes in docs/paper-gaps/.
% Notation follows the LDT paper (references/ldt-paper/).

\newcommand{\Lean}{\textsc{Lean}}
\newcommand{\leanid}[1]{\textsf{#1}}
\newcommand{\ghissue}[1]{\href{https://github.com/LionSR/MIPStarRE/issues/#1}{GitHub issue \##1}}
\newcommand{\ghpr}[1]{\href{https://github.com/LionSR/MIPStarRE/pull/#1}{GitHub PR \##1}}

% --- Paper-compatible notation ---
\newcommand{\F}{\mathbb{F}}
\newcommand{\N}{\mathbb{N}}
\newcommand{\C}{\mathbb{C}}
\newcommand{\tr}{\operatorname{tr}}
\newcommand{\ot}{\otimes}
\newcommand{\E}{\mathop{\bf E\/}}
\newcommand{\bx}{\mathbf{x}}
\newcommand{\by}{\mathbf{y}}
\newcommand{\bu}{\mathbf{u}}
\newcommand{\bv}{\mathbf{v}}

% Approximation relation (paper uses \approx_\eps and \simeq_\zeta)
\newcommand{\approxeps}{\approx_{\varepsilon}}
\newcommand{\simgeq}{\simeq_{\zeta}}

% Polynomial spaces (paper uses \polyfunc, \polysub, \polymeas)
\newcommand{\polyfunc}[3]{\operatorname{Poly}(#1,#2,#3)}
\newcommand{\polysub}[3]{\operatorname{PolySub}(#1,#2,#3)}
\newcommand{\polymeas}[3]{\operatorname{PolyMeas}(#1,#2,#3)}


\title{Issue 1230: How the Self-Improvement SDP Is Used}
\author{}
\date{2026-05-05}

\begin{document}
\maketitle

\section{At a glance}

\begin{itemize}
  \item \textbf{Difficulty: hard.}  The issue is the full
  \textbf{strong-duality and complementary-slackness producer} for the
  self-improvement SDP, not just the already-fixed submeasurement interface.
  \item \textbf{Estimated weight:} about \textbf{8--12 theorem interfaces} and
  \textbf{1000--2000 lines} if developed project-locally.  The estimate includes
  canonicalization, Slater hypotheses, value equality, slackness, saturation, and
  the self-improvement adapter.
  \item \textbf{Mathlib/project split:} roughly \textbf{35\% Mathlib} and
  \textbf{65\% project-local}.  Mathlib supplies matrix order, traces,
  positive-semidefinite cones, convex cones, and separation tools; the project
  must build the paper-specific block SDP and helper interface.
  \item \textbf{Key Mathlib inputs:} positive semidefinite matrices, trace pairing,
  convex cones and dual cones, finite-dimensional linear maps, and hyperplane
  separation for closed/proper cones.%
  \footnote{Representative APIs live around
  \texttt{Mathlib.LinearAlgebra.Matrix.PosDef},
  \texttt{Mathlib.Analysis.Matrix.Order},
  \texttt{Mathlib.Analysis.Convex.Cone.Dual},
  \texttt{Mathlib.Analysis.Convex.Cone.InnerDual}, and
  \texttt{Mathlib.Geometry.Convex.Cone.*}.}
\end{itemize}

\section{Key theorem forms to develop}

The following theorem forms are the ones needed for the LDT self-improvement
argument.  They are deliberately stronger than the local issue~196
submeasurement fix.

\begin{enumerate}
  \item \textbf{PSD effect bounds.}  For every polynomial $g$, prove
  \[
    0\preceq A_g\preceq I,
    \qquad
    A_g=\E_{\bu} A^{\bu}_{g(\bu)}.
  \]
  This supplies dual feasibility of the strict witness $Z=2I$ and nonnegative
  objective gain when slack mass is completed.

  \item \textbf{Canonical block SDP equivalence.}  The paper primal over
  submeasurements should be equivalent to the canonical PSD variable
  $X=\operatorname{diag}(T_{g_1},\ldots,T_{g_M},S)$ with
  \[
    \sum_gT_g+S=I,
    \qquad
    \langle C,X\rangle=\sum_g\tr(T_gA_g).
  \]
  Its dual should be exactly the paper dual $Z\succeq A_g$.

  \item \textbf{Slater and strong duality.}  The strict witnesses
  $T_g=(2M)^{-1}I$ and $Z=2I$ should trigger finite-dimensional SDP strong
  duality: there are optimal $X_*$ and $Z_*$ with equal primal and dual values.

  \item \textbf{Saturated slackness witness.}  Choose an optimal primal witness
  with \textbf{no slack block}, i.e. $\sum_gT_g=I$, and a dual optimum $Z$ such
  that
  \[
    Z\succeq A_g,
    \qquad
    T_g(Z-A_g)=0
    \quad\text{for all }g.
  \]
  Saturation can be obtained by completing primal slack into any outcome without
  decreasing the objective; slackness then follows from zero duality gap and
  positive-semidefinite dual slacks.

  \item \textbf{Self-improvement SDP output.}  Package the previous item as the
  exact data consumed by the helper proof:
  \[
    T\text{ is a measurement},
    \qquad
    Z\succeq0,
    \qquad
    Z\succeq A_g,
    \qquad
    T_gZ=T_gA_g.
  \]
  If the current Lean interface still asks for the extra field $I\preceq Z$,
  either prove it for the selected witness or, more faithfully, relax that
  interface to the paper's $Z\succeq0$ and $Z\succeq A_g$ fields.%
  \footnote{The current bridge isolating this interface issue is
  \doclink{MIPStarRE/LDT/SelfImprovement/Theorems/Results/SdpMatrixBridge.lean}.}
\end{enumerate}

\section{The source SDP}

The self-improvement lemma in \cite{Ji2020LowIndividualDegree} uses one SDP.%
\footnote{The SDP is in \path{references/ldt-paper/self_improvement.tex},
lines~62--88.  The strict-feasibility and complementary-slackness paragraph is
in lines~168--190.}
For $A_g=\E_{\bu}A^{\bu}_{g(\bu)}$, the primal and dual are
\begin{align}
  \sup_{\{T_g\}}&\quad \sum_g\tr(T_gA_g),&
  T_g&\succeq0,& \sum_gT_g&\preceq I,\tag{\texttt{eq:primal-objective}}\\
  \inf_Z&\quad \tr(Z),&
  Z&\succeq A_g&&\text{for every }g.\tag{\texttt{eq:dual-objective}}
\end{align}
The primal object is a \textbf{submeasurement}.  The SDP lemma then selects an
optimal pair with
\[
  \sum_gT_g=I,
  \qquad
  T_gZ=T_gA_g.
  \tag{\texttt{lem:sdp}, \texttt{eq:slater}}
\]

\section{How the proof uses the optimum}

The paper defines
\[
  H^u_h=A^u_{h(u)}T_hA^u_{h(u)},
  \qquad
  H_h=\E_{\bu}H^{\bu}_h.
  \tag{\texttt{self\_improvement.tex:203--220}}
\]
The proof uses \textbf{exactly three SDP outputs}: submeasurement bounds make
$H^u$ and $H$ submeasurements; saturation collapses sums over $T_g$; and
slackness replaces $A_g$ by $Z$ after left multiplication by $T_g$.  For example,
\[
  \E_{\bu}\sum_h\bra{\psi}(T_hA^{\bu}_{h(\bu)})\ot I\ket{\psi}
  =\sum_h\bra{\psi}(T_hZ)\ot I\ket{\psi}.
\]
Dual feasibility $Z\succeq A_g$ then supplies the boundedness certificate carried
through the induction.%
\footnote{The two slackness substitutions are in
\path{references/ldt-paper/self_improvement.tex}, lines~397--400 and 580--584.
The boundedness conclusion is restated in
\path{references/ldt-paper/inductive_step.tex}, lines~277--284, 315--322,
477--484, and 542--549.}

\section{Blueprint and formalization}

The blueprint matches the paper: submeasurement primal, strict feasible witness,
then a saturated slackness-producing optimum.%
\footnote{See \path{blueprint/src/chapter/ch07_self_improvement.tex}, especially
\texttt{lem:sdp-uniform-feasible-witness},
\texttt{lem:sdp-matrix-feasible-bounds},
\texttt{lem:sdp-matrix-slackness-output}, and \texttt{lem:sdp}.}
Lean already keeps the submeasurement boundary explicit; the remaining gap is
the \textbf{producer} for the stronger slackness-carrying statement.%
\footnote{Representative declarations are \leanid{sdpStrictPrimalSubMeas},
\leanid{sdpPrimalObjective}, \leanid{sdpComplementarySlacknessEquation}, and
\leanid{averagedSandwichedPolynomialSubMeas} in
\doclink{MIPStarRE/LDT/SelfImprovement/Defs.lean}.  The interfaces are
\leanid{SdpStatement} and \leanid{SdpStatementWithSlackness} in
\doclink{MIPStarRE/LDT/SelfImprovement/Theorems/Statements.lean}.}

\section{Conclusion}

There is \textbf{no new discrepancy} in the paper's SDP use.  The hard remaining
work is to prove or import strong duality with complementary slackness, specialize
it to the paper's block SDP, produce a saturated optimal witness, and connect it
to the self-improvement interface.

\bibliographystyle{alpha}
\bibliography{references}

\end{document}
